\chapter{Stock Effects and Draw Parameters}
\label{ch:stock-effects}

XNA shipped five built-in shader effects, and CNA ports all five. This chapter covers the
shared \cnaclass{Effect} base and its supporting types, then each stock effect, and finally the
recurring cross-renderer defect patterns exposed by their verification.

Many histories below come from the project's original EasyGL/Vulkan/BGFX comparison campaign.
``All three'' means those three measured implementations, not all \RendererFamilyCount{} current
renderer families;
closed defects are retained as mutation-tested evidence, not presented as present limitations.

\section{Effect: stock and compiled modes}

\cnaclass{Effect} inherits \cnaclass{GraphicsResource} and is not copyable. The
\texttt{Effect(device)} constructor creates the small stock-effect base graph. The byte-buffer
constructor now accepts a bounded, structurally validated XNA/FNA Direct3D~9 Effect Framework
binary --- the \texttt{.fxb}-shaped payload normally stored by XNA's general
\texttt{EffectReader}. It rejects an empty or over-64-MiB buffer, malformed container graphs, and
the distinct MonoGame \texttt{MGFX} container. It then requires
\texttt{GraphicsCapability::CompiledEffects}; a renderer that did not opt in throws
\texttt{NotSupportedException} before any draw.

FNA3D supports this route in its normal build. SDL\_GPU, the EasyGL profiles, and Vulkan support it
only when their respective \texttt{CNA\_*\_COMPILED\_EFFECTS} option is enabled; those options are
off by default because they add MojoShader. Other renderers report the capability false. This is
compiled Effect Framework support, not a compiler for \texttt{.fx} or HLSL source, not DXBC or
MGFX ingestion, and not the renderer-specific \cnaclass{ShaderEffect} API. Chapter
\ref{ch:shader-fx-gap} separates the formats and backends.

\texttt{Clone()} is virtual and, as a
\texttt{CNAEXT} deviation from FNA's garbage-collected reference semantics, returns an owning
raw pointer rather than a managed reference --- there is no GC in C\texttt{++} to hand a
shared reference back to. \texttt{Apply()} itself is \texttt{CNAEXT} --- real FNA has no
public \texttt{Effect.Apply()}, only \texttt{EffectPass.Apply()} --- and the protected virtual
\texttt{FillGpuDrawParams(GpuDrawParams\&)} is the single mechanism every stock effect below
uses to hand its own parameters down to whichever renderer is active; the base implementation
is a no-op, and every stock effect overrides it.

\cnaclass{EffectParameter} carries a name, semantic, row/column counts, class/type enums, and a
broad family of typed \texttt{GetValue*} / \texttt{SetValue*} accessors. In a compiled effect,
those records are reflected from the binary and mutable values synchronize to the renderer
runtime before a pass is applied. In a stock effect, the collection remains the narrower manual
shadow described in \S\ref{sec:stock-effect-parameter-boundary}. The
parameter, pass, and technique collections all store elements behind \texttt{unique\_ptr} rather
than by value, so a reference obtained earlier --- including a captured \texttt{CurrentTechnique}
--- survives a later \texttt{Add()} reallocation. \cnaclass{EffectPass::Apply()} throws
\texttt{InvalidOperationException} if an owned pass does not belong to the effect's currently
selected technique, including when \texttt{CurrentTechnique} is null --- matching FNA's own
``Applied a pass not in the current technique!'' guard, mapped onto a defined C\texttt{++}
exception rather than the null-reference crash the equivalent null case produces in C\#.

\section{Stock parameters are not compiled-effect reflection}
\label{sec:stock-effect-parameter-boundary}

The XNA-looking \texttt{Parameters} property needs a mode-qualified reading. A compiled effect
reflects parameters, techniques, passes, annotations, arrays and structures from its binary;
current technique and mutable values are copied by \texttt{Clone()}, and pass application pushes
the renderer's blend, depth/stencil, rasterizer, sampler, texture and shader state changes through
the device. A bare stock-mode \cnaclass{Effect}, by contrast, starts with one manually created
\texttt{Default} technique containing one \texttt{P0} pass and with \emph{zero} parameters.

The five stock effects do not all even present the same structural approximation.
\cnaclass{BasicEffect} adds no parameter records at all, even though real FNA's compiled
BasicEffect exposes its parsed parameter table:

\begin{lstlisting}
auto& parameters = basic.getParametersProperty();
EffectParameter* diffuse = parameters["DiffuseColor"]; // nullptr
\end{lstlisting}

\cnaclass{AlphaTestEffect},
\cnaclass{DualTextureEffect}, \cnaclass{EnvironmentMapEffect}, and \cnaclass{SkinnedEffect}
each call a private \texttt{CacheEffectParameters()} helper instead. It manually creates a
small set of internal shadows: respectively 6, 5, 12, and 12 records for the derived material,
fog, matrix, and shader-index calculations. Those lists are neither a bytecode reflection result
nor a full public description of the effect; for example, the manually cached AlphaTest and
DualTexture lists contain no texture parameter despite each effect having a texture property.

The normal render flow is consequently one-way:

\begin{lstlisting}
stockEffect.setDiffuseColorProperty(tint); // authoritative effect field / dirty bit
for (EffectPass& pass : stockEffect.getCurrentTechniqueProperty()->getPassesProperty()) {
    pass.Apply();                           // OnApply may refresh internal shadow records
    graphicsDevice.DrawIndexedPrimitives(/* ... */);
    // GraphicsDevice calls stockEffect.FillGpuDrawParams(), which reads the effect fields.
}
\end{lstlisting}

Writing a manually found \cnaclass{EffectParameter} therefore does not form a portable
alternative to a stock effect's named property and does not create an arbitrary renderer uniform.
The device dispatches \texttt{FillGpuDrawParams()} directly, not the collection. A few
properties deliberately use a shadow as their own storage --- fog color is the important example
in the four effects that cache it --- but that is an implementation detail that a later property
assignment or \texttt{OnApply()} calculation can overwrite. Use the public
\texttt{setDiffuseColorProperty}, \texttt{setTextureProperty}, matrix, lighting, fog, and
skinning APIs for stock material state; use the explicit \texttt{ShaderEffect::SetUniform*}
surface, with Chapter~\ref{ch:shader-fx-gap}'s renderer limits, for a custom uniform.

\subsection{What an individual parameter retains}

An \cnaclass{EffectParameter} stores independent C\texttt{++} float, integer, string, and
texture-pointer caches. Its metadata does not check that a requested setter matches the declared
class, type, shape, or count; it is descriptive data, not runtime validation. Array getters
return at most the currently stored entries rather than necessarily the requested count. A
negative request is caller-invalid: vector/matrix/quaternion and Boolean getters happen to yield
an empty result, while the integer and float array implementations form an invalid iterator range.
No focused test covers this mismatch or negative-count path.

The texture overloads deserve the same caution. CNA keeps a separate raw slot for each texture
overload. Calling \texttt{SetValue(Texture*)} therefore does not populate the matching typed
texture getter, even when the dynamic object has that derived type; there is no generic texture
getter. FNA has one texture reference which its typed getters cast. CNA also accepts and returns
strings where FNA's \texttt{SetValue(string)} explicitly throws
\texttt{NotImplementedException}. None of these pointers retains a texture, and none is a GPU
binding instruction. The standalone parameter tests cover same-overload scalar/vector/matrix/
quaternion/texture round trips and defaults, but not metadata mismatch, truncation, negative
counts, generic-texture behavior, lifetime, stock-effect linkage, or rendering.

\subsection{Pass selection still matters}

\cnaclass{Effect::Apply()} is a useful CNA-only convenience: it calls \texttt{OnApply()} and
makes the effect current on the device even if \texttt{CurrentTechnique} is null or names an
unrelated technique. It does not itself select or validate a pass. The familiar loop above is
safe because \cnaclass{EffectPass::Apply()} compares its owning technique's stable identity with
the current one before reaching that convenience method. The public C\texttt{++} constructors
also allow an ownerless pass; applying such a pass simply returns without setting any effect,
where FNA constructs passes internally. Keep application code on the current technique's own
passes and do not use null/unrelated techniques or ownerless manually assembled passes as a
configuration API.

\section{GpuDrawParams: the shared stock-effect draw packet}

Every stock effect's \texttt{FillGpuDrawParams} override populates the same struct, defined
once in the renderer contract layer (Chapter~\ref{ch:backend-architecture}) and forwarded from
\texttt{Effect::FillGpuDrawParams()} to whichever renderer is currently active. It carries two
ordinary texture slots plus a dedicated environment-map cube slot, diffuse/ambient/emissive/%
specular colors, three directional lights (each with direction, diffuse, and specular),
the world matrix, alpha-test parameters, Fresnel/environment-map parameters, a 72-bone
skinning palette with a weights-per-vertex count, fog parameters, a set of mode flags
(\texttt{textureEnabled}, \texttt{vertexColorEnabled}, \texttt{lightingEnabled},
\texttt{dualTexture}, \texttt{envMapping}, \texttt{skinned}, \texttt{pbr}), instancing fields,
a pointer for a custom \texttt{ShaderEffect}-compiled program, and six PBR-specific
metallic-roughness fields. Two flags are worth flagging directly: \texttt{preferPerPixelLighting}
and \texttt{specularEnabled}. Earlier project documentation called both flags ``ignored by
every renderer except D3D9.'' That statement is historical; the current scoped result appears
after \cnaclass{SkinnedEffect} below.

\section{BasicEffect}

\cnaclass{BasicEffect} is the largest of the five: public \texttt{World} / \texttt{View}/%
\texttt{Projection} fields, \texttt{VertexColorEnabled}, the full material color set
(\texttt{DiffuseColor}, \texttt{EmissiveColor}, \texttt{SpecularColor}, \texttt{SpecularPower},
\texttt{Alpha}, \texttt{AmbientLightColor}), \texttt{LightingEnabled}/%
\texttt{PreferPerPixelLighting}, \texttt{TextureEnabled} / \texttt{Texture}, full
\texttt{DirectionalLight0/1/2}, and \texttt{EnableDefaultLighting()}. Its 22-property audit
found and fixed several early default-value bugs (\texttt{VertexColorEnabled},
\texttt{DirectionalLight0.Enabled}, and \texttt{DirectionalLight.Direction}'s default value
were all initially wrong), then a sequence of real, per-renderer rendering bugs: the
no-texture \texttt{VertexColorEnabled} toggle was ignored by all three of EasyGL, Vulkan, and
BGFX (three separate bugs, one per renderer); the textured-plus-vertex-color-plus-diffuse case
dropped \texttt{DiffuseColor} entirely on EasyGL and BGFX (Vulkan was already correct); and a
shared bug across all three renderers never checked \texttt{DirectionalLight0.Enabled} at all,
so a disabled light kept lighting the surface regardless. BGFX carried one additional,
much wider bug: its vertex-layout builder never declared \texttt{Normal} / \texttt{TexCoord0}
attributes for any stride except the skinned 52-byte layout --- silently breaking lit-normal
\emph{and} UV interpolation for every other stride, invisible until then because every prior
test happened to use a UV-insensitive $1\times1$ texture. \texttt{EmissiveColor} was dropped
entirely in the no-lighting path on all three renderers; ambient/emissive/specular forwarding
and specular highlights were corrected afterward. The focused composition fixture then produced
byte-identical output on EasyGL, Vulkan, and BGFX at the pinned revisions; this result does not
generalize to untested effect combinations or renderer families.

\subsection{Lit, textured cube}

\cnaclass{BasicEffect} is deliberately the effect a first 3D CNA program reaches for, and its
own real headers confirm why it reads like plain XNA rather than a wrapped shader system:
\texttt{World} and \texttt{VertexColorEnabled} are ordinary public fields (matching real XNA's
own field-not-property choice for this specific effect), while \texttt{Texture} and
\texttt{TextureEnabled} go through the usual \texttt{getXProperty} / \texttt{setXProperty} pair:

\begin{lstlisting}
BasicEffect effect(getGraphicsDeviceProperty());
effect.World = Matrix::CreateRotationY(rotationAngle);
effect.setViewProperty(camera.View());
effect.setProjectionProperty(camera.Projection());

effect.setTextureEnabledProperty(true);
effect.setTextureProperty(&cubeTexture);

effect.EnableDefaultLighting();          // sets DirectionalLight0/1/2 to XNA's own defaults
effect.DirectionalLight0.setDiffuseColorProperty(Vector3(1.0f, 0.95f, 0.9f)); // warm key light

for (EffectPass& pass : effect.getCurrentTechniqueProperty()->getPassesProperty()) {
    pass.Apply();
    getGraphicsDeviceProperty().DrawIndexedPrimitives(/* ... */);
}
\end{lstlisting}

\texttt{EnableDefaultLighting()} is worth calling before touching any individual
\texttt{DirectionalLightN} property, not after --- reading its own implementation directly
confirms it unconditionally overwrites \texttt{DirectionalLight0}'s diffuse color, direction,
specular color, and enabled flag (and likewise for lights 1 and 2) every time it runs; setting
\texttt{DirectionalLight0}'s diffuse color first and calling \texttt{EnableDefaultLighting()}
second would silently overwrite the custom color right back to the default. This is exactly the kind of ordering trap the per-renderer bug list below
does \emph{not} cover, because it is a real API-usage hazard rather than a renderer bug ---
\texttt{EnableDefaultLighting()}'s own real behavior, faithfully ported, not a CNA-specific
quirk.

\section{AlphaTestEffect}

\cnaclass{AlphaTestEffect} implements \texttt{IEffectMatrices} and \texttt{IEffectFog} only
--- no lighting. \texttt{AlphaFunction} (a \cnaclass{CompareFunction}, defaulting to
\texttt{Greater}) and \texttt{ReferenceAlpha} (an unclamped 0--255 int) implement the actual
alpha test; \texttt{VertexColorEnabled}, \texttt{DiffuseColor} / \texttt{Alpha}, and
\texttt{Texture} round out the surface. Its audit found zero property-default bugs (the
first test coverage for this effect written from scratch), and a full 8-value
\cnaclass{CompareFunction} sweep passed 24 of 24 cases on the three campaign renderers. The
campaign initially found that \texttt{VertexColorEnabled} had no effect on Vulkan or BGFX because
their alpha-test pipeline declared no color attribute. The correction added
stride-24 sibling vertex shaders and discriminating pixel tests on both products; EasyGL was
already correct. Fog forwarding was a real, fixed bug on EasyGL
specifically for this effect --- and, in the course of fixing it, revealed a much larger
finding: fog was a \emph{total} no-op on both Vulkan and BGFX for every 3D effect, because no
shader file on either renderer mentioned fog at all. That project-wide gap was closed in one
pass, adding real fog uniforms, varyings, and blend formula to every 3D shader on both
renderers at once --- the same ``one shared fix closes it everywhere'' pattern already seen
for \cnaclass{SpriteBatch}'s flip bug in Chapter~\ref{ch:spritebatch}. A separate, unrelated
BGFX-only bug --- seven texture-binding call sites with no fallback-to-white-texture
else-branch --- was found and fixed at the same time, affecting every effect that binds a
possibly-null texture, not just this one.

\subsection{Cutout-alpha chain-link fence}

The classic real use of \cnaclass{AlphaTestEffect} is a cutout texture --- foliage, a
chain-link fence, a wire mesh --- where a pixel is either fully opaque or fully discarded, with
no blended edge at all (unlike ordinary alpha \emph{blending}, which this effect deliberately
does not do). \texttt{AlphaFunction} and \texttt{ReferenceAlpha} together express the test as a
comparison against the texture's own alpha channel:

\begin{lstlisting}
AlphaTestEffect effect(getGraphicsDeviceProperty());
effect.setWorldProperty(fenceWorld);
effect.setViewProperty(camera.View());
effect.setProjectionProperty(camera.Projection());

effect.setTextureProperty(&fenceTexture);         // alpha channel: 255 = wire, 0 = gap
effect.setAlphaFunctionProperty(CompareFunction::Greater);
effect.setReferenceAlphaProperty(128);            // keep only pixels with alpha > 128

for (EffectPass& pass : effect.getCurrentTechniqueProperty()->getPassesProperty()) {
    pass.Apply();
    getGraphicsDeviceProperty().DrawIndexedPrimitives(/* ... */);
}
\end{lstlisting}

The corresponding regression protects this shape. A stride-24
\cnaclass{VertexPositionColorTexture} fence must multiply its per-vertex color only when
\texttt{VertexColorEnabled} is true, and the combined alpha must participate in the discard.
Reverting the Vulkan/BGFX sibling shaders reproduced the old diffuse-only value and even changed
the pass/discard decision; restoring them returned both cases to the expected pixels.

\section{DualTextureEffect}

Two independent texture slots (\texttt{Texture}, \texttt{Texture2}), plus
\texttt{VertexColorEnabled}, \texttt{DiffuseColor} / \texttt{Alpha}, and fog --- no lighting.
The audit's headline finding was a shared formula bug across all three campaign renderers:
FNA's real shader doubles the first texture's RGB value (\texttt{color.rgb *= 2}) before
multiplying it against the overlay texture, and none of the three shaders initially
implemented that factor --- invisible to every prior test, because a test built from
0-or-1-saturated color values cannot distinguish ``multiplied by one'' from ``multiplied by
two and then saturated.'' The factor is present in all three pinned shader paths now.
\texttt{Texture2}'s null-texture fallback had its own,
BGFX-specific bug (again, no else-branch), fixed alongside the project-wide fallback fix
above. Fog needed its own separate EasyGL fix here, because this effect uses its own
dedicated shader with no fog infrastructure at all, unlike \cnaclass{AlphaTestEffect}'s
shared shader path; Vulkan and BGFX were already covered by the project-wide fog fix.
\texttt{VertexColorEnabled} was another closed defect: the correction added a stride-24 colored
variant on EasyGL, Vulkan, and BGFX and verified enabled and disabled cases independently.

\subsection{Lightmapped floor}

\cnaclass{DualTextureEffect}'s two independent texture slots are the classic pre-baked-lighting
pattern from the original Xbox 360/Windows XNA era: a base color texture plus a second,
low-resolution lightmap multiplied over it, with no runtime lighting computation needed at all:

\begin{lstlisting}
DualTextureEffect effect(getGraphicsDeviceProperty());
effect.setWorldProperty(floorWorld);
effect.setViewProperty(camera.View());
effect.setProjectionProperty(camera.Projection());

effect.setTextureProperty(&floorDiffuse);   // base color, tiled across the floor mesh's own UVs
effect.setTexture2Property(&floorLightmap); // low-res baked lighting, its own second UV channel

for (EffectPass& pass : effect.getCurrentTechniqueProperty()->getPassesProperty()) {
    pass.Apply();
    getGraphicsDeviceProperty().DrawIndexedPrimitives(/* ... */);
}
\end{lstlisting}

The audit's headline finding directly changes what result this produces: real FNA's shader
doubles \texttt{Texture}'s RGB value before multiplying it against \texttt{Texture2} ---
\texttt{color.rgb *= 2} --- specifically so a lightmap value of 0.5 (mid-gray, the ``neutral,
unlit'' value most lightmap bakers emit) reproduces the base texture's own color exactly, rather
than darkening it to half brightness. A lightmap authored against real XNA/FNA's doubling
behavior will render visibly too dark on any CNA build predating the fix that added this factor
to all three campaign shaders at once. A colored lightmapped floor uses the stride-24
variant; the pinned EasyGL, Vulkan, and BGFX routes now distinguish the enabled and disabled
\texttt{VertexColorEnabled} cases.

\section{EnvironmentMapEffect}

\cnaclass{EnvironmentMapEffect} combines \texttt{IEffectMatrices}, \texttt{IEffectLights}, and
\texttt{IEffectFog}. It accepts a diffuse \texttt{Texture}, a cube
\texttt{EnvironmentMap}, \texttt{EnvironmentMapAmount},
\texttt{EnvironmentMapSpecular}, \texttt{FresnelFactor}, \texttt{EmissiveColor}, and lighting
state. As in FNA, \texttt{LightingEnabled} is always true and its setter throws when passed
\texttt{false}. The 14-property API audit found no property-contract defects.

\begin{historynote}
The stock-effect audit corrected four rendering/state defects. \texttt{Clone()} omitted
\texttt{FogColor} here and in \cnaclass{AlphaTestEffect}, \cnaclass{DualTextureEffect}, and
\cnaclass{SkinnedEffect}. The cube-map blend was additive instead of FNA's \texttt{lerp}; its
base and specular terms omitted the cube alpha scale; and all renderers lacked Fresnel
edge-weighting. CNA evaluates the corrected Fresnel term per pixel, whereas FNA evaluates it per
vertex. EasyGL and BGFX also used the world matrix for normals under non-uniform scale; they now
use the required transpose-inverse, as Vulkan already did.
\end{historynote}

A capstone fixture composes these corrections and produces byte-identical output on EasyGL,
Vulkan, and BGFX. That result is scoped to the fixture and those renderer routes.

\subsection{Chrome car body}

A reflective, Fresnel-edge-brightened surface --- the archetypal use for this effect --- needs
both the diffuse texture and the \cnaclass{TextureCube} environment map bound, plus lighting
enabled so the Fresnel and specular terms below have a light direction to work from:

\begin{lstlisting}
EnvironmentMapEffect effect(getGraphicsDeviceProperty());
effect.setWorldProperty(carBodyWorld);
effect.setViewProperty(camera.View());
effect.setProjectionProperty(camera.Projection());

effect.setTextureProperty(&carBodyPaint);
effect.setEnvironmentMapProperty(&skybox);        // the previously created TextureCube
effect.setEnvironmentMapAmountProperty(0.85f);    // mostly reflective, not a full mirror
effect.setFresnelFactorProperty(1.0f);
effect.setLightingEnabledProperty(true);          // required -- setter throws if set false

for (EffectPass& pass : effect.getCurrentTechniqueProperty()->getPassesProperty()) {
    pass.Apply();
    getGraphicsDeviceProperty().DrawIndexedPrimitives(/* ... */);
}
\end{lstlisting}

At the pin, \texttt{FresnelFactor} brightens reflections toward grazing angles and
\texttt{EnvironmentMapAmount} interpolates between the base color and reflection. Earlier
implementations omitted the Fresnel term and used an additive cube-map blend, which could
over-brighten the result. EasyGL and BGFX also once transformed normals by the world matrix;
their current transpose-inverse transform preserves reflection shape under non-uniform scale.

\subsection{Three-light forwarding and a stale plan entry}

\begin{historynote}
At the pin, \texttt{plan\_graphics.md} still marks Task~890---forwarding
\texttt{DirectionalLight1} and \texttt{DirectionalLight2} through
\cnaclass{EnvironmentMapEffect}---as open. The implementation has moved ahead of that checkbox.
\texttt{FillGpuDrawParams()} copies both lights' direction and diffuse fields when each light is
enabled, and the renderer implementations consume those fields at their uniform-update sites.
The EasyGL source comment retains the task number, which makes the stale plan entry traceable.
FNA's one-light shader variant is also selected when lights 1 and 2 are disabled.
\end{historynote}

A three-light scene distinguishes the current forwarding path from the former one-light result;
a single-light scene cannot:

\begin{lstlisting}
EnvironmentMapEffect effect(getGraphicsDeviceProperty());
effect.setWorldProperty(carBodyWorld);
effect.setViewProperty(camera.View());
effect.setProjectionProperty(camera.Projection());
effect.setTextureProperty(&carBodyPaint);
effect.setEnvironmentMapProperty(&skybox);
effect.setLightingEnabledProperty(true);

effect.DirectionalLight0.setEnabledProperty(true);           // key light, warm
effect.DirectionalLight0.setDiffuseColorProperty(Vector3(1.0f, 0.85f, 0.7f));
effect.DirectionalLight1.setEnabledProperty(true);            // fill light, cool
effect.DirectionalLight1.setDiffuseColorProperty(Vector3(0.2f, 0.3f, 0.5f));
effect.DirectionalLight2.setEnabledProperty(true);            // rim light, from behind
effect.DirectionalLight2.setDiffuseColorProperty(Vector3(0.6f, 0.6f, 0.6f));

for (EffectPass& pass : effect.getCurrentTechniqueProperty()->getPassesProperty()) {
    pass.Apply();
    getGraphicsDeviceProperty().DrawIndexedPrimitives(/* ... */);
}
\end{lstlisting}

Before the fix, this scene ignored \texttt{DirectionalLight1} and
\texttt{DirectionalLight2}. At the pinned revision all three enabled lights reach the renderer.

\section{SkinnedEffect}

\cnaclass{SkinnedEffect} accepts at most \texttt{MaxBones = 72}. Its
\texttt{WeightsPerVertex} property accepts only 1, 2, or 4 and throws for other values.
\texttt{SetBoneTransforms} / \texttt{GetBoneTransforms(count)}, specular lighting, fog, and
three directional lights follow the stock-effect route. CNA adds
\texttt{VertexColorEnabled} for imported glTF meshes with a \texttt{COLOR\_0} attribute; XNA's
\cnaclass{SkinnedEffect} has no corresponding property.

Identity, single-bone, and two-bone weighted skinning are pixel-verified on the three campaign
renderers. Historical defects affected surrounding state: \texttt{Clone()} dropped
\texttt{SpecularColor} and \texttt{SpecularPower}; lights 1 and 2 were not forwarded; the
specular fields had no GPU implementation; and \texttt{WeightsPerVertex} did not gate the
four-weight sum. The pinned implementation preserves the cloned fields, forwards all enabled
lights, evaluates a Blinn--Phong half-vector specular term, and gates the sum for the selected
weight count.

\subsection{Driving SkinnedEffect from AnimationPlayer}

Chapter~\ref{ch:models-meshes}'s own worked \cnaclass{AnimationPlayer} example ended with
exactly the call this effect exists to receive --- \texttt{SetBoneTransforms} takes
\texttt{GetSkinTransforms()}'s output directly, already pre-multiplied by each bone's inverse
bind pose, which is specifically the form the GPU skinning shader expects:

\begin{lstlisting}
SkinnedEffect effect(getGraphicsDeviceProperty());
effect.setViewProperty(camera.View());
effect.setProjectionProperty(camera.Projection());
effect.setWeightsPerVertexProperty(2);    // must be 1, 2, or 4 -- throws otherwise
effect.setSpecularColorProperty(Vector3(0.3f, 0.3f, 0.3f));
effect.setSpecularPowerProperty(16.0f);
effect.EnableDefaultLighting();

// Every frame, after player.Update(...):
effect.SetBoneTransforms(player.GetSkinTransforms());
\end{lstlisting}

Rigs exceeding 72 palette joints require mesh partitioning into draws with disjoint palettes.
Set \texttt{WeightsPerVertex} to match the asset. Projects that previously requested two weights
but depended on the former four-weight behavior will render differently at the pin.

\section{Closing the \texttt{preferPerPixelLighting} divergence}

\begin{contractnote}
\texttt{PreferPerPixelLighting} defaults to \texttt{false}, matching XNA's per-vertex
(Gouraud-interpolated) default. The pinned \texttt{D3D9}, \texttt{EasyGL}, \texttt{Vulkan},
\texttt{BGFX}, \texttt{D3D11}, and \texttt{D3D12} paths honor the choice. \texttt{WebGPU}
honors it for \cnaclass{BasicEffect}; \cnaclass{SkinnedEffect} and
\cnaclass{EnvironmentMapEffect} are not implemented there. It is inapplicable to the 2D-only
\texttt{SdlRenderer} and non-rasterizing \texttt{Headless} renderer. The \texttt{Software}
rasterizer has no lighting engine, so this remains an open limitation there. Do not extrapolate
this campaign result to unlisted implementation families.
\end{contractnote}

Earlier shared dispatch omitted both \texttt{PreferPerPixelLighting} and
\cnaclass{EnvironmentMapEffect}'s \texttt{specularEnabled}. The affected GPU paths therefore
used fragment lighting regardless of the public setting. D3D9 already contained Microsoft's
per-vertex and per-pixel shader families; the other covered renderers added per-vertex variants
using FNA's lighting formula in the vertex stage. \cnaclass{AlphaTestEffect} and
\cnaclass{DualTextureEffect} have no lighting path and were unchanged.

The new discriminating fixtures exposed four secondary issues:

\begin{longtable}{@{}p{0.15\textwidth}p{0.73\textwidth}@{}}
\toprule
Renderer & Finding and disposition \\
\midrule
\endhead
\texttt{D3D9} & Pixel-lit constants were uploaded only to the vertex stage, leaving the pixel
shader's lighting registers zero. The upload now targets the stage that evaluates the light. \\
\texttt{Vulkan} & New pipeline-cache maps were absent from teardown. A validation-layer report
at \texttt{vkDestroyDevice()} exposed the leaked \texttt{VkPipeline} objects. \\
\texttt{WebGPU} & Center values 187/203 differed from linear 127/152 because the preferred
swapchain format applies sRGB encoding; the computed encodings match the observations. \\
\texttt{D3D12} & A specular-exactness fixture depended on per-fragment evaluation. It now sets
\texttt{PreferPerPixelLighting} to \texttt{true}, preserving that test's intended geometry. \\
\bottomrule
\end{longtable}

Lit-scene baselines changed when the XNA-compatible default became effective: low-poly surfaces
can appear more faceted and highlights can soften under vertex interpolation. Those changes are
expected for the default path; the D3D12 fixture above explicitly selects the alternate path.

\section{The pattern behind these bugs}

The audits found three recurring failure modes:

\begin{itemize}
  \item a missing shader term, such as doubled dual-texture RGB, Fresnel weighting, or specular;
  \item a public flag that never reached GPU state, such as the former
        \texttt{VertexColorEnabled} and \texttt{WeightsPerVertex} gaps; and
  \item a shared dispatch defect affecting several renderers, including the former light-enable,
        fog-forwarding, and clone-state omissions.
\end{itemize}

The tested core interpolation and lighting formulas were not broadly defective; the observed
failures were specific omissions. This distinction keeps the evidence scoped to the paths and
fixtures that exposed them.
